% RASCUNHO OPCIONAL — NAO PUBLICAR.
% O pipeline global descrito aqui nao possui execucao completa. Este arquivo
% contem marcadores @@...@@ e nao fornece os resultados dos artigos atuais.
% Para a submissao corrigida, use exclusivamente: overlief/IEEE/artigo.tex.
\documentclass[10pt,conference,a4paper]{IEEEtran}
\IEEEoverridecommandlockouts

\usepackage[utf8]{inputenc}
\usepackage[T1]{fontenc}
\usepackage[brazil]{babel}
\usepackage{amsmath,amssymb}
\usepackage{array}
\usepackage{booktabs}
\usepackage{graphicx}
\usepackage{microtype}
\usepackage{tikz}
\usepackage{url}
\usetikzlibrary{arrows.meta,calc,fit,positioning}

\graphicspath{{../figuras/}}
\emergencystretch=2em
\renewcommand{\abstractname}{Resumo}
\renewcommand{\IEEEkeywordsname}{Palavras-chave}
\newcounter{algoritmo}

\title{Avaliação do DeepNPTS e de Modelos de Referência para Previsão Mensal
de Irradiância Global Horizontal em Localidades de Fábricas de Veículos
Elétricos}

\author{Gustavo de Melo Nascimento, Thales Wulfert Cabral e Gustavo Fraidenraich%
\thanks{Os autores são da Faculdade de Engenharia Elétrica e de Computação,
Universidade Estadual de Campinas, Campinas, SP, Brasil. E-mails:
g239679@dac.unicamp.br, t264377@dac.unicamp.br e gfraiden@unicamp.br.}}

\begin{document}

\maketitle

\begin{abstract}
Este artigo avalia o Deep Non-Parametric Time Series Forecaster (DeepNPTS) na
previsão de séries mensais de Irradiância Global Horizontal (GHI), com
horizonte de um mês. A avaliação considera dez séries relacionadas, cujas
coordenadas correspondem a localidades de fábricas de veículos elétricos, e
dados da National Solar Radiation Database entre 2019 e 2024. As séries são
agregadas por mês civil, normalizadas por min--max e quantizadas em 128 níveis
com parâmetros estimados apenas no treinamento. Uma única rede global do
DeepNPTS discreto, proveniente da implementação oficial do GluonTS, é comparada
à persistência, ao sazonal ingênuo, à climatologia mensal, ao XGBoost, ao
perceptron multicamadas, à rede neural recorrente, à LSTM, à DilatedRNN e ao
DeepAR. O protocolo
retrospectivo emprega 48 alvos de treinamento e 12 origens de teste por
localidade, cinco sementes e previsão 
\textit{walk-forward} de um passo sem reajuste. O menor MAE médio entre
localidades foi obtido por \textbf{@@MODELO-MELHOR@@}, com
\textbf{@@MAE-MELHOR@@ W/m$^2$}. O DeepNPTS alcançou MAE de
\textbf{@@MAE-DEEPNPTS@@ W/m$^2$} e posição \textbf{@@POSICAO-DEEPNPTS@@} no
ranqueamento. Para DeepNPTS e DeepAR, CRPS, cobertura e largura dos intervalos
de 90\% também são analisados. Os resultados caracterizam exclusivamente o
protocolo e a janela retrospectiva estudados.
\end{abstract}

\begin{IEEEkeywords}
Irradiância solar, GHI, previsão probabilística, séries temporais, DeepNPTS,
DeepAR.
\end{IEEEkeywords}

\section{Introdução}
\label{sec:introducao}

A disponibilidade do recurso solar condiciona a produção de sistemas
fotovoltaicos e influencia decisões de dimensionamento, planejamento e
operação energética. A previsão de irradiância contribui para reduzir a
incerteza associada à geração solar; sua formulação, porém, depende do horizonte
temporal, da resolução e das informações meteorológicas disponíveis
\cite{antonanzas2016,voyant2017,yang2018}.

Este trabalho concentra-se na Irradiância Global Horizontal (GHI), isto é, na
irradiância total que incide sobre uma superfície horizontal. Em termos das
componentes direta e difusa,

\begin{equation}
  \mathrm{GHI}=\mathrm{DHI}+\mathrm{DNI}\cos(\theta_z),
  \label{eq:ghi}
\end{equation}

\noindent em que DHI é a irradiância difusa horizontal, DNI é a irradiância
direta normal e $\theta_z$ é o ângulo zenital solar \cite{duffie2013}. A GHI é
empregada por expressar, em W/m$^2$, a disponibilidade solar sobre o plano
horizontal e por integrar produtos padronizados da National Solar Radiation
Database (NSRDB) \cite{sengupta2018}. Neste estudo, a variável-alvo é a média
mensal de GHI, não energia em Wh/m$^2$ nem produção fotovoltaica.

A previsão de uma série mensal com horizonte de um mês pode apoiar a comparação
de localidades e análises preliminares de planejamento. A agregação mensal
atenua flutuações intradiárias e explicita a sazonalidade, mas reduz o número de
amostras a doze por ano \cite{antonanzas2016,voyant2022}. Portanto, modelos
flexíveis precisam ser confrontados com referências sazonais simples, sobretudo
quando o histórico contém poucos ciclos anuais.

A literatura de previsão solar abrange métodos estatísticos, aprendizado de
máquina e redes neurais \cite{voyant2017,sobri2018}. Neste trabalho, os modelos
pontuais aprendidos incluem XGBoost \cite{chen2016}, Perceptron Multicamadas
(MLP) \cite{hornik1989}, Rede Neural Recorrente (RNN) \cite{elman1990}, Long
Short-Term Memory (LSTM) \cite{hochreiter1997} e DilatedRNN
\cite{chang2017}. Persistência, sazonal ingênuo e climatologia mensal compõem
as referências determinísticas.

Avanços em previsão global permitem ajustar um único modelo a várias séries
relacionadas. O DeepAR, por exemplo, aprende uma distribuição preditiva por
meio de uma rede recorrente autorregressiva global \cite{salinas2020}. O Deep
Non-Parametric Time Series Forecaster (DeepNPTS) também explora séries
relacionadas, mas constrói a distribuição preditiva a partir de valores do
próprio contexto histórico, com probabilidades aprendidas por uma rede global
\cite{rangapuram2023}. Essa propriedade é pertinente a séries curtas porque
restringe a variante discreta a valores já representados no contexto, embora
não garanta melhor acurácia em qualquer aplicação.

O objetivo deste artigo é avaliar a implementação discreta oficial do
DeepNPTS, disponível no GluonTS \cite{alexandrov2020}, para prever a GHI média
do mês seguinte em dez localidades, comparando-a sob o mesmo protocolo temporal
com seis modelos aprendidos e três referências simples. As fábricas definem
apenas os pontos geográficos das consultas; não são utilizados dados de carga,
produção industrial, geração fotovoltaica ou consumo de energia dessas
instalações.

\subsection{Principais contribuições}

As contribuições deste estudo são:

\begin{enumerate}
  \item avaliação do DeepNPTS discreto como um único modelo global treinado nas
  dez séries relacionadas, com amostragem probabilística e identidade da
  localidade como atributo categórico estático;
  \item comparação temporal homogênea com três referências simples e seis
  modelos aprendidos, incluindo arquiteturas tabulares, recorrentes e
  probabilísticas;
  \item protocolo auditável com transformações ajustadas apenas no treinamento,
  validação cronológica, cinco sementes, previsões linha a linha e métricas
  pontuais e probabilísticas; e
  \item análise explícita das limitações de uma janela retrospectiva curta, sem
  alegação de pioneirismo, superioridade universal ou validação prospectiva.
\end{enumerate}

O restante do artigo está organizado da seguinte forma. A
Seção~\ref{sec:sistema} descreve o sistema e o DeepNPTS. A
Seção~\ref{sec:configuracao} apresenta dados, protocolo, modelos e métricas. A
Seção~\ref{sec:resultados} discute os resultados, e a
Seção~\ref{sec:conclusoes} apresenta as conclusões.

\section{Sistema de previsão mensal de GHI}
\label{sec:sistema}

A Fig.~\ref{fig:sistema} apresenta o fluxo completo. Os blocos são explicados
na mesma ordem da figura e formalizados no Algoritmo~\ref{alg:sistema}.

\begin{figure*}[!t]
\centering
\resizebox{0.98\textwidth}{!}{%
\begin{tikzpicture}[
  grupo/.style={draw,dashed,rounded corners=1pt,minimum height=3.5cm,inner sep=4pt},
  caixa/.style={draw,align=center,minimum height=0.72cm,font=\scriptsize,fill=white},
  titulo/.style={align=center,font=\scriptsize\bfseries},
  rotulo/.style={font=\scriptsize},
  seta/.style={-{Latex[length=2.1mm]},thick}
]
\node[grupo,minimum width=2.45cm] (ga) at (0,0) {};
\node[titulo] at ([yshift=-0.33cm]ga.north) {Dez séries\\relacionadas};
\node[caixa,minimum width=1.75cm] (dados) at ([yshift=-1.70cm]ga.north)
  {GHI diária\\NSRDB\\2019--2024};
\node[rotulo] at ([yshift=-0.27cm]ga.south) {(a)};

\node[grupo,minimum width=2.65cm] (gb) at (2.95,0) {};
\node[titulo] at ([yshift=-0.33cm]gb.north) {Auditoria e\\agregação};
\node[caixa,minimum width=2.05cm] (audit) at ([yshift=-1.28cm]gb.north)
  {Cobertura, lacunas,\\duplicidades e sinais};
\node[caixa,minimum width=2.05cm,below=0.25cm of audit] (mensal)
  {Média por\\mês civil};
\draw[seta] (audit)--(mensal);
\node[rotulo] at ([yshift=-0.27cm]gb.south) {(b)};

\node[grupo,minimum width=3.05cm] (gc) at (6.15,0) {};
\node[titulo] at ([yshift=-0.33cm]gc.north) {Transformação e\\representações};
\node[caixa,minimum width=2.35cm] (transforma) at ([yshift=-1.18cm]gc.north)
  {Min--max do treino\\+ quantização};
\node[caixa,minimum width=2.35cm,below=0.23cm of transforma] (janelas)
  {Janelas causais,\\atributos e categoria};
\draw[seta] (transforma)--(janelas);
\node[rotulo] at ([yshift=-0.27cm]gc.south) {(c)};

\node[grupo,minimum width=3.45cm] (gd) at (9.95,0) {};
\node[titulo] at ([yshift=-0.33cm]gd.north) {Treinamento global\\e referências};
\node[caixa,minimum width=2.75cm] (aprendidos) at ([yshift=-1.17cm]gd.north)
  {XGBoost, MLP, RNN, LSTM,\\DilatedRNN, DeepAR, DeepNPTS};
\node[caixa,minimum width=2.75cm,below=0.23cm of aprendidos] (baselines)
  {Persistência, sazonal\\e climatologia};
\node[rotulo] at ([yshift=-0.27cm]gd.south) {(d)};

\node[grupo,minimum width=2.95cm] (ge) at (13.65,0) {};
\node[titulo] at ([yshift=-0.33cm]ge.north) {Teste retrospectivo\\walk-forward};
\node[caixa,minimum width=2.25cm] (origens) at ([yshift=-1.17cm]ge.north)
  {Um mês à frente\\sem reajuste};
\node[caixa,minimum width=2.25cm,below=0.23cm of origens] (amostras)
  {Ponto, amostras\\e intervalos};
\draw[seta] (origens)--(amostras);
\node[rotulo] at ([yshift=-0.27cm]ge.south) {(e)};

\node[grupo,minimum width=2.75cm] (gf) at (17.00,0) {};
\node[titulo] at ([yshift=-0.33cm]gf.north) {Avaliação};
\node[caixa,minimum width=2.05cm] (metrica) at ([yshift=-1.18cm]gf.north)
  {Métricas por\\localidade};
\node[caixa,minimum width=2.05cm,below=0.23cm of metrica] (macro)
  {Macro-MAE e\\incerteza};
\draw[seta] (metrica)--(macro);
\node[rotulo] at ([yshift=-0.27cm]gf.south) {(f)};

\draw[seta] (dados.east)--(audit.west);
\draw[seta] (mensal.east)--(transforma.west);
\draw[seta] (janelas.east)--(aprendidos.west);
\draw[seta] (aprendidos.east)--(origens.west);
\draw[seta] (amostras.east)--(metrica.west);
\end{tikzpicture}%
}
\caption{Visão geral do sistema: (a) séries de GHI; (b) auditoria e
agregação; (c) transformação e construção das entradas; (d) ajuste dos modelos
e cálculo das referências; (e) previsão retrospectiva de um passo; e
(f) avaliação por localidade e agregação macro.}
\label{fig:sistema}
\end{figure*}

\noindent\textbf{Estágio (a) -- séries relacionadas.} Cada localidade
$\ell\in\mathcal L$ fornece uma série diária $S_\ell$. As dez séries possuem a
mesma variável e frequência, mas escalas e sazonalidades locais distintas. A
identidade $\ell$ é preservada em todas as previsões.

\noindent\textbf{Estágio (b) -- auditoria e agregação.} O sistema rejeita datas
inválidas, duplicidades diárias, lacunas de cobertura, valores não finitos e
GHI negativa. Em seguida, calcula a média de cada mês civil. Os arquivos
preservados contêm médias diárias derivadas da consulta de 60 minutos; as
respostas horárias brutas não são reconstituídas.

\noindent\textbf{Estágio (c) -- transformação e entradas.} Para cada
localidade, os limites min--max são estimados sem usar os alvos de teste. Os
valores transformados são discretizados uniformemente. Modelos tabulares
recebem defasagens consecutivas, médias móveis, calendário e codificação da
localidade. Redes recorrentes recebem a sequência cronológica por um ramo e as
covariáveis por outro. DeepNPTS e DeepAR recebem diretamente as séries
normalizadas e a categoria estática da localidade.

\noindent\textbf{Estágio (d) -- treinamento e referências.} Cada modelo
aprendido é global: seus parâmetros são compartilhados pelas dez séries. A
seleção de épocas ou árvores, quando aplicável, utiliza uma janela cronológica
interna, e o modelo final é reajustado em todo o treinamento. As três
referências não possuem parâmetros aprendidos conjuntamente.

\noindent\textbf{Estágio (e) -- teste retrospectivo.} Em cada origem, o
histórico é truncado no último mês disponível e o mês seguinte é previsto. O
valor de referência de um mês de 2024 torna-se disponível para a origem
seguinte, mas os parâmetros dos modelos permanecem fixos. Portanto, o protocolo
não equivale a prever conjuntamente os doze meses de 2024 em dezembro de 2023.

\noindent\textbf{Estágio (f) -- avaliação.} Todas as previsões retornam à escala
física em W/m$^2$. Calculam-se métricas em cada localidade e, depois, a média
não ponderada entre localidades. DeepNPTS e DeepAR também produzem amostras e
intervalos preditivos.

\subsection{DeepNPTS global discreto}

A Fig.~\ref{fig:deepnpts} detalha o modelo central do estudo. Um único
DeepNPTS é ajustado com instâncias extraídas das dez séries. Para uma janela de
contexto $\mathbf z_{\ell,t}=(z_{\ell,t-C+1},\ldots,z_{\ell,t})$, a rede
feedforward combina valores passados, indicadores de observação, atributos de
tempo e o \textit{embedding} da localidade. Sua saída define probabilidades
$p_{\ell,t,j}$ sobre as $C$ posições do contexto.

\begin{figure}[!t]
\centering
\resizebox{\columnwidth}{!}{%
\begin{tikzpicture}[
  caixa/.style={draw,rounded corners=1pt,align=center,minimum height=0.83cm,
                minimum width=1.85cm,font=\scriptsize},
  seta/.style={-{Latex[length=1.9mm]},thick},
  node distance=0.34cm
]
\node[caixa] (series) {Dez séries\\de treinamento};
\node[caixa,right=of series] (instancias) {Contextos $\mathbf z_{\ell,t}$,\\tempo e localidade};
\node[caixa,right=of instancias,fill=blue!5] (rede) {Rede global\\compartilhada};
\node[caixa,right=of rede] (prob) {Probabilidades\\$p_{\ell,t,1:C}$};
\node[caixa,right=of prob,fill=green!8] (dist) {Distribuição discreta\\$\sum_jp_j\delta_{z_j}$};
\node[caixa,below=0.42cm of dist] (saida) {Amostras, mediana\\e quantis};
\draw[seta] (series)--(instancias);
\draw[seta] (instancias)--(rede);
\draw[seta] (rede)--(prob);
\draw[seta] (prob)--(dist);
\draw[seta] (dist)--(saida);
\end{tikzpicture}%
}
\caption{DeepNPTS global discreto. A rede aprende, nas séries relacionadas,
probabilidades sobre os valores que compõem o contexto de cada previsão.}
\label{fig:deepnpts}
\end{figure}

A distribuição preditiva de um passo é

\begin{equation}
  \widehat F_{\ell,t+1}(x)
  =\sum_{j=1}^{C}p_{\ell,t,j}\,
  \mathbb{I}\!\left(z_{\ell,t-C+j}\leq x\right),
  \quad \sum_{j=1}^{C}p_{\ell,t,j}=1,
  \label{eq:deepnpts_dist}
\end{equation}

\noindent em que $\mathbb{I}(\cdot)$ é a função indicadora. Equivalentemente,
$J\sim\mathrm{Categorical}(p_{\ell,t,1:C})$ e
$\widetilde z_{\ell,t+1}=z_{\ell,t-C+J}$. A repetição da amostragem representa
a incerteza da previsão. A variante avaliada é a rede discreta oficial do
GluonTS 0.16.2; não se trata do suavizador determinístico por similaridade e
recência utilizado em experimentos anteriores do projeto.

O ponto associado ao MAE é a mediana da mistura das amostras produzidas nas
sementes. Para os demais modelos aprendidos pontuais, utiliza-se a média das
previsões por semente:

\begin{equation}
 \widehat y_i^{\mathrm{ens}}=\frac{1}{|\mathcal S|}
 \sum_{s\in\mathcal S}\widehat y_{i,s}.
 \label{eq:ensemble}
\end{equation}

\subsection{Procedimento experimental}

O Algoritmo~\ref{alg:sistema} usa apenas símbolos; os valores adotados são
apresentados na Seção~\ref{sec:configuracao}.

\begin{figure*}[!t]
\centering
\refstepcounter{algoritmo}\label{alg:sistema}
\textbf{Algoritmo \thealgoritmo: Treinamento global, previsão e avaliação}

\fbox{%
\begin{minipage}{0.94\textwidth}
\footnotesize
\begin{tabbing}
\quad\=\quad\=\kill
\textbf{Entrada:} séries $\mathcal D$, localidades $\mathcal L$, modelos $\mathcal M$,
sementes $\mathcal S$, contexto $C$, horizonte $h$, corte $c$, níveis $Q$ e origens $\mathcal O$.\\
\textbf{Para cada} $\ell\in\mathcal L$ \textbf{faça}\\
\> $S_\ell\leftarrow$ auditar e agregar $\mathcal D_\ell$ por mês civil;\\
\> $T_\ell\leftarrow$ ajustar min--max antes de $c$ e quantizar em $Q$ níveis;\\
\> $(D_\ell^{tr},D_\ell^{va},D_\ell^{te})\leftarrow$
construir janelas causais de contexto $C$ e horizonte $h$;\\
\textbf{fim}\\
\textbf{Para cada} $m\in\mathcal M$ \textbf{e} $s\in\mathcal S$ \textbf{faça}\\
\> $\eta_m\leftarrow$ fixar a configuração; quando aplicável, selecionar sua complexidade em $D^{tr}$ e $D^{va}$;\\
\> $\theta_{m,s}\leftarrow$ ajustar ou reajustar $m$ no conjunto pré-teste completo com $\eta_m$;\\
\> \textbf{Para cada} origem $o\in\mathcal O$ \textbf{faça}\\
\>\> $H_o\leftarrow$ truncar cada série na origem $o$;\\
\>\> $\widehat F_{m,s,o}\leftarrow$ prever $h$ passos com $\theta_{m,s}$ e $H_o$;\\
\> \textbf{fim}\\
\textbf{fim}\\
$\widehat F_m\leftarrow$ consolidar sementes; calcular métricas por localidade e médias macro;\\
$m^\star\leftarrow\arg\min_{m\in\mathcal M}\overline{\mathrm{MAE}}_m$;\\
\textbf{Saída:} previsões, amostras, intervalos, métricas, $m^\star$ e manifesto.
\end{tabbing}
\end{minipage}}
\end{figure*}

\section{Configuração experimental e métricas}
\label{sec:configuracao}

\subsection{Base de dados e protocolo temporal}

Foram utilizados dados do produto GOES Aggregated PSM v4 da NSRDB
\cite{sengupta2018,nlr2026}. A NSRDB combina informação de satélite e modelos
físicos; seus valores são tratados como referência modelada, e não como
medição de solo em sentido metrológico. A consulta original adotou intervalo de
60 minutos. O projeto preserva médias das 24 horas de cada dia, inclusive o
período noturno, que são posteriormente agregadas pela média de cada mês civil.

As localidades são BYD Camaçari, BMW San Luis Potosí, Ford Rouge Electric
Vehicle Center, GM Factory Zero, Hyundai Metaplant Georgia, Lucid AMP 1 Casa
Grande, Rivian Normal, Tesla Fremont Factory, Tesla Gigafactory Nevada e Tesla
Gigafactory Texas. Cada arquivo contém 2.192 dias de 1 de janeiro de 2019 a 31
de dezembro de 2024; a auditoria dos dez arquivos totalizou 21.920 linhas e não
encontrou datas inválidas, duplicidades, lacunas diárias, GHI não finita ou GHI
negativa.

Com contexto de 12 meses, cada localidade fornece 60 alvos: 48 no treinamento,
de janeiro de 2020 a dezembro de 2023, e 12 no teste, de janeiro a dezembro de
2024. O ano de 2019 fornece o contexto dos primeiros exemplos. A janela de 2024
foi inspecionada durante o desenvolvimento do projeto; assim, a análise é
retrospectiva e exploratória, não uma confirmação prospectiva em dados nunca
vistos.

A Tabela~\ref{tab:setup} resume o protocolo. A normalização e a quantização são
ajustadas separadamente por localidade. Na seleção cronológica interna, os
limites são estimados somente antes da validação; no reajuste final, usam-se
apenas dados até dezembro de 2023. Os modelos aprendidos recebem valores
quantizados, enquanto as referências simples operam nos valores físicos exatos.
Essa diferença faz parte do protocolo declarado e deve ser considerada na
interpretação.

\begin{table}[!t]
\centering
\caption{Configuração do experimento mensal.}
\label{tab:setup}
\small
\setlength{\tabcolsep}{3.5pt}
\renewcommand{\arraystretch}{1.05}
\begin{tabular}{>{\raggedright\arraybackslash}p{2.05cm}>{\raggedright\arraybackslash}p{5.95cm}}
\toprule
\textbf{Item} & \textbf{Configuração} \\
\midrule
Fonte e variável & NSRDB, GOES Aggregated PSM v4; média mensal de GHI em W/m$^2$. \\
Cobertura & 2019--2024; dez localidades; 72 meses por série. \\
Coleta preservada & Médias diárias derivadas de consultas com intervalo de 60 min; sem respostas horárias brutas. \\
Contexto e horizonte & 12 meses de contexto; previsão de um mês. \\
Particionamento & 48 alvos de treino (2020--2023) e 12 alvos de teste (2024) por localidade. \\
Teste & 12 origens \textit{walk-forward}, valor anterior disponível e modelo fixo. \\
Transformação & Min--max ajustado somente no pré-teste e quantização uniforme em 128 níveis. \\
Sementes & 11, 23, 42, 67 e 89. \\
Amostragem & 500 amostras por semente e origem para DeepNPTS/DeepAR; 2.500 na mistura final. \\
Hardware & @@HARDWARE-CPU-RAM-SO@@. \\
Software & Python 3.12.1; GluonTS 0.16.2; TensorFlow 2.21; PyTorch 2.6; versões completas no manifesto. \\
\bottomrule
\end{tabular}
\end{table}

\subsection{Modelos e hiperparâmetros}

As entradas tabulares possuem 12 defasagens consecutivas, médias móveis de 3,
6 e 12 meses, seno e cosseno do mês-alvo e codificação \textit{one-hot} da
localidade. RNN, LSTM e DilatedRNN recebem a sequência $12\times1$ e, em um ramo
auxiliar, médias móveis, calendário e localidade. A DilatedRNN implementa
conexões recorrentes de salto $h_t=\phi(Wx_t+Uh_{t-s}+b)$ com $s\in\{1,2,4\}$,
em vez de tratar colunas tabulares como passos temporais.

DeepNPTS e DeepAR são os estimadores PyTorch oficiais do GluonTS 0.16.2. Ambos
compartilham parâmetros entre as dez séries e usam a localidade como categoria
estática. O DeepNPTS discreto utiliza duas camadas ocultas de 12 unidades e
\textit{embedding} de dimensão 5, conforme os padrões determinados pela
cardinalidade na versão fixada da biblioteca. A Tabela~\ref{tab:hiper} apresenta
os demais hiperparâmetros.

\begin{table*}[!t]
\centering
\caption{Hiperparâmetros dos modelos aprendidos.}
\label{tab:hiper}
\small
\setlength{\tabcolsep}{4pt}
\renewcommand{\arraystretch}{1.06}
\begin{tabular}{>{\raggedright\arraybackslash}p{2.2cm}>{\raggedright\arraybackslash}p{14.8cm}}
\toprule
\textbf{Modelo} & \textbf{Configuração} \\
\midrule
XGBoost & Profundidade 3; taxa 0,03; peso mínimo do filho 4; \textit{subsample}=0,8; \textit{colsample}=0,8; $L_1=0,1$; $L_2=10$; até 1.000 árvores, parada após 40 rodadas sem melhora em 2023 e reajuste com a quantidade escolhida. \\
MLP & Camadas 64 e 32; ReLU; otimizador L-BFGS; $\alpha=0,05$; máximo de 2.000 iterações; padronização interna das entradas. \\
RNN & SimpleRNN com 16 unidades \textit{tanh}; ramo denso de 16 unidades ReLU; saída linear; Adam 0,001; MSE; lote 32; até 300 épocas e paciência 30. \\
LSTM & LSTM com 16 unidades \textit{tanh}; ramo denso de 16 unidades ReLU; saída linear; Adam 0,001; MSE; lote 32; até 300 épocas e paciência 30. \\
DilatedRNN & Três camadas recorrentes com 16 unidades e dilatações 1, 2 e 4; ramo denso 16 ReLU; saída linear; Adam 0,001; MSE; lote 32; até 300 épocas e paciência 30. \\
DeepNPTS & Variante discreta global; contexto 12; camadas 12--12; categoria estática; log-verossimilhança discreta; Adam $10^{-5}$; 100 épocas; 100 lotes/época; lote 32; sem escalonamento interno; 500 amostras por semente e origem. \\
DeepAR & Duas camadas recorrentes de 40 unidades; saída Student-$t$ e log-verossimilhança; \textit{dropout}=0,1; lags 1--12; categoria estática; Adam 0,001; 100 épocas; 50 lotes/época; lote 32; sem escalonamento interno; 500 amostras por semente e origem. \\
\bottomrule
\end{tabular}
\end{table*}

A quantidade de árvores e de épocas é escolhida usando 2023 como validação
cronológica, sem o teste. Depois, XGBoost, RNN, LSTM e DilatedRNN são
reinicializados e ajustados em 2020--2023 com a complexidade selecionada.
DeepNPTS e DeepAR usam a configuração fixa da Tabela~\ref{tab:hiper}. O
processo é repetido para cinco sementes. Todos os ajustes são executados em CPU.

\subsection{Métricas pontuais e seleção}

Sejam $y_{\ell,i}$ a GHI de referência da NSRDB,
$\widehat y_{\ell,i,m}$ a previsão do modelo $m$, $\bar y_\ell$ a média dos
$n_\ell$ alvos de teste e $e_{\ell,i,m}=y_{\ell,i}-\widehat y_{\ell,i,m}$. As
métricas pontuais são

\begin{align}
 \mathrm{MAE}_{\ell,m}
 &=\frac{1}{n_\ell}\sum_{i=1}^{n_\ell}|e_{\ell,i,m}|,
 \label{eq:mae}\\
 \mathrm{MSE}_{\ell,m}
 &=\frac{1}{n_\ell}\sum_{i=1}^{n_\ell}e_{\ell,i,m}^{2},
 \label{eq:mse}\\
 \mathrm{RMSE}_{\ell,m}
 &=\sqrt{\frac{1}{n_\ell}\sum_{i=1}^{n_\ell}e_{\ell,i,m}^{2}},
 \label{eq:rmse}\\
 R^2_{\ell,m}
 &=1-\frac{\sum_{i=1}^{n_\ell}e_{\ell,i,m}^{2}}
 {\sum_{i=1}^{n_\ell}(y_{\ell,i}-\bar y_\ell)^2},
 \label{eq:r2}\\
 \mathrm{nRMSE}_{\ell,m}(\%)
 &=100\,\frac{\mathrm{RMSE}_{\ell,m}}{\bar y_\ell}.
 \label{eq:nrmse}
\end{align}

MAE e RMSE retornam a W/m$^2$; MSE é expresso em
W$^2$/m$^4$; $R^2$ é adimensional; e nRMSE expressa erro relativo à média da
localidade. MAE é a métrica primária por sua interpretação direta e menor
sensibilidade a desvios extremos que métricas quadráticas. As demais métricas
descrevem propriedades complementares do erro \cite{hyndman2006}.

O critério primário é a média não ponderada do MAE nas localidades:

\begin{equation}
 \overline{\mathrm{MAE}}_m=\frac{1}{|\mathcal L|}
 \sum_{\ell\in\mathcal L}\mathrm{MAE}_{\ell,m},
 \qquad
 m^\star=\arg\min_{m\in\mathcal M}\overline{\mathrm{MAE}}_m.
 \label{eq:macro}
\end{equation}

O rótulo $m^\star$ identifica apenas o menor erro neste protocolo, não um
modelo universalmente superior. Intervalos de 95\% para o macro-MAE são obtidos
por \textit{bootstrap} de localidades. Diferenças pareadas
$\Delta_{\ell,m}=\mathrm{MAE}_{\ell,m}-\mathrm{MAE}_{\ell,\mathrm{clim}}$ são
comparadas com zero pelo teste de Wilcoxon, com correção de Holm para múltiplas
comparações; valores positivos favorecem a climatologia.

\subsection{Métricas probabilísticas}

Para DeepNPTS e DeepAR, sejam $x_{i,b}$ a $b$-ésima de $B$ amostras
preditivas e $q_{i,u}$ o quantil de nível $u$. O CRPS empírico é
\cite{gneiting2007}

\begin{equation}
 \mathrm{CRPS}_i=
 \frac{1}{B}\sum_{b=1}^{B}|x_{i,b}-y_i|
 -\frac{1}{2B^2}\sum_{b=1}^{B}\sum_{c=1}^{B}|x_{i,b}-x_{i,c}|.
 \label{eq:crps}
\end{equation}

Para o intervalo preditivo central de nível $1-\alpha$, calculam-se

\begin{align}
 \mathrm{PICP}
 &=\frac{1}{n}\sum_{i=1}^{n}
 \mathbb{I}\!\left(q_{i,\alpha/2}\le y_i\le q_{i,1-\alpha/2}\right),
 \label{eq:picp}\\
 \mathrm{MPIW}
 &=\frac{1}{n}\sum_{i=1}^{n}
 \left(q_{i,1-\alpha/2}-q_{i,\alpha/2}\right).
 \label{eq:mpiw}
\end{align}

CRPS mede simultaneamente calibração e concentração e é menor quando a
distribuição está mais próxima do alvo. PICP mede a cobertura empírica, e MPIW,
em W/m$^2$, mede a largura média. Cobertura isolada não basta, pois intervalos
arbitrariamente largos podem cobrir mais observações \cite{gneiting2007}.

\section{Resultados e discussão}
\label{sec:resultados}

\subsection{Comparação média dos modelos}

A Tabela~\ref{tab:metricas_medias} apresenta as métricas da previsão
consolidada, ordenadas pelo macro-MAE. O desvio de sementes se refere ao
macro-MAE das cinco execuções; não representa a dispersão entre localidades.

\begin{table*}[!t]
\centering
\caption{Métricas médias entre as dez localidades no teste retrospectivo.}
\label{tab:metricas_medias}
\small
\setlength{\tabcolsep}{5pt}
\renewcommand{\arraystretch}{1.04}
\begin{tabular}{lrrrrrr}
\toprule
\textbf{Modelo} & \textbf{MAE} & \textbf{DP seeds} & \textbf{MSE} & \textbf{RMSE} & \textbf{$R^2$} & \textbf{nRMSE} \\
 & \textbf{W/m$^2$} & \textbf{W/m$^2$} & \textbf{W$^2$/m$^4$} & \textbf{W/m$^2$} & & \textbf{\%} \\
\midrule
Persistência     & @@MAE-PERSISTENCIA@@ & -- & @@MSE-PERSISTENCIA@@ & @@RMSE-PERSISTENCIA@@ & @@R2-PERSISTENCIA@@ & @@NRMSE-PERSISTENCIA@@ \\
Sazonal ingênuo  & @@MAE-SAZONAL@@ & -- & @@MSE-SAZONAL@@ & @@RMSE-SAZONAL@@ & @@R2-SAZONAL@@ & @@NRMSE-SAZONAL@@ \\
Climatologia     & @@MAE-CLIMATOLOGIA@@ & -- & @@MSE-CLIMATOLOGIA@@ & @@RMSE-CLIMATOLOGIA@@ & @@R2-CLIMATOLOGIA@@ & @@NRMSE-CLIMATOLOGIA@@ \\
XGBoost          & @@MAE-XGBOOST@@ & @@DP-XGBOOST@@ & @@MSE-XGBOOST@@ & @@RMSE-XGBOOST@@ & @@R2-XGBOOST@@ & @@NRMSE-XGBOOST@@ \\
MLP              & @@MAE-MLP@@ & @@DP-MLP@@ & @@MSE-MLP@@ & @@RMSE-MLP@@ & @@R2-MLP@@ & @@NRMSE-MLP@@ \\
RNN              & @@MAE-RNN@@ & @@DP-RNN@@ & @@MSE-RNN@@ & @@RMSE-RNN@@ & @@R2-RNN@@ & @@NRMSE-RNN@@ \\
LSTM             & @@MAE-LSTM@@ & @@DP-LSTM@@ & @@MSE-LSTM@@ & @@RMSE-LSTM@@ & @@R2-LSTM@@ & @@NRMSE-LSTM@@ \\
DilatedRNN       & @@MAE-DILATEDRNN@@ & @@DP-DILATEDRNN@@ & @@MSE-DILATEDRNN@@ & @@RMSE-DILATEDRNN@@ & @@R2-DILATEDRNN@@ & @@NRMSE-DILATEDRNN@@ \\
DeepAR           & @@MAE-DEEPAR@@ & @@DP-DEEPAR@@ & @@MSE-DEEPAR@@ & @@RMSE-DEEPAR@@ & @@R2-DEEPAR@@ & @@NRMSE-DEEPAR@@ \\
DeepNPTS         & @@MAE-DEEPNPTS@@ & @@DP-DEEPNPTS@@ & @@MSE-DEEPNPTS@@ & @@RMSE-DEEPNPTS@@ & @@R2-DEEPNPTS@@ & @@NRMSE-DEEPNPTS@@ \\
\bottomrule
\end{tabular}
\end{table*}

\begin{figure}[!t]
\centering
\IfFileExists{../figuras/mae_medio_modelos.pdf}{%
  \includegraphics[width=\columnwidth]{mae_medio_modelos.pdf}%
}{\fbox{\parbox[c][3.7cm][c]{0.92\columnwidth}{\centering
Figura pendente: \texttt{mae\_medio\_modelos.pdf}}}}
\caption{Macro-MAE dos modelos. As barras de erro representam o desvio entre
sementes quando aplicável; valores menores são melhores.}
\label{fig:mae_medio}
\end{figure}

No protocolo avaliado, \textbf{@@MODELO-MELHOR@@} apresentou o menor
macro-MAE, de \textbf{@@MAE-MELHOR@@ W/m$^2$}, com intervalo de 95\%
\textbf{[@@IC-INF-MELHOR@@; @@IC-SUP-MELHOR@@]}. O DeepNPTS ficou na posição
\textbf{@@POSICAO-DEEPNPTS@@}, com diferença de
\textbf{@@DELTA-DEEPNPTS-CLIM@@ W/m$^2$} em relação à climatologia; o intervalo
pareado foi \textbf{[@@IC-INF-DEEPNPTS-CLIM@@;
@@IC-SUP-DEEPNPTS-CLIM@@]} e o valor $p$ ajustado por Holm foi
\textbf{@@P-HOLM-DEEPNPTS-CLIM@@}. Esses valores
\textbf{@@INTERPRETACAO-ESTATISTICA-DEEPNPTS@@}.

O desempenho relativo deve ser interpretado à luz de apenas quatro ciclos
anuais para ajuste dos alvos. Referências sazonais podem ser competitivas nesse
regime, enquanto redes mais parametrizadas precisam estimar relações globais a
partir de poucas instâncias temporais. Por outro lado, compartilhar parâmetros
entre localidades aumenta o conjunto de instâncias, mas não cria novos anos de
variabilidade climática.

\subsection{Resultados do DeepNPTS por localidade}

A Tabela~\ref{tab:deepnpts_local} detalha as métricas do DeepNPTS. O menor MAE
ocorreu em \textbf{@@LOCAL-MELHOR-DEEPNPTS@@}
(\textbf{@@MAE-LOCAL-MELHOR-DEEPNPTS@@ W/m$^2$}), e o maior em
\textbf{@@LOCAL-PIOR-DEEPNPTS@@}
(\textbf{@@MAE-LOCAL-PIOR-DEEPNPTS@@ W/m$^2$}). A amplitude foi
\textbf{@@AMPLITUDE-MAE-DEEPNPTS@@ W/m$^2$}. Essa variação mostra que a média
global não descreve uniformemente as dez localidades.

\begin{table*}[!t]
\centering
\caption{Desempenho pontual do DeepNPTS por localidade.}
\label{tab:deepnpts_local}
\small
\setlength{\tabcolsep}{5pt}
\renewcommand{\arraystretch}{1.03}
\begin{tabular}{lrrrrr}
\toprule
\textbf{Localidade} & \textbf{MAE} & \textbf{MSE} & \textbf{RMSE} & \textbf{$R^2$} & \textbf{nRMSE} \\
 & \textbf{W/m$^2$} & \textbf{W$^2$/m$^4$} & \textbf{W/m$^2$} & & \textbf{\%} \\
\midrule
BMW San Luis Potosí & @@BMW-MAE@@ & @@BMW-MSE@@ & @@BMW-RMSE@@ & @@BMW-R2@@ & @@BMW-NRMSE@@ \\
BYD Camaçari & @@BYD-MAE@@ & @@BYD-MSE@@ & @@BYD-RMSE@@ & @@BYD-R2@@ & @@BYD-NRMSE@@ \\
Ford Rouge & @@FORD-MAE@@ & @@FORD-MSE@@ & @@FORD-RMSE@@ & @@FORD-R2@@ & @@FORD-NRMSE@@ \\
GM Factory Zero & @@GM-MAE@@ & @@GM-MSE@@ & @@GM-RMSE@@ & @@GM-R2@@ & @@GM-NRMSE@@ \\
Hyundai Georgia & @@HYUNDAI-MAE@@ & @@HYUNDAI-MSE@@ & @@HYUNDAI-RMSE@@ & @@HYUNDAI-R2@@ & @@HYUNDAI-NRMSE@@ \\
Lucid Casa Grande & @@LUCID-MAE@@ & @@LUCID-MSE@@ & @@LUCID-RMSE@@ & @@LUCID-R2@@ & @@LUCID-NRMSE@@ \\
Rivian Normal & @@RIVIAN-MAE@@ & @@RIVIAN-MSE@@ & @@RIVIAN-RMSE@@ & @@RIVIAN-R2@@ & @@RIVIAN-NRMSE@@ \\
Tesla Fremont & @@FREMONT-MAE@@ & @@FREMONT-MSE@@ & @@FREMONT-RMSE@@ & @@FREMONT-R2@@ & @@FREMONT-NRMSE@@ \\
Tesla Nevada & @@NEVADA-MAE@@ & @@NEVADA-MSE@@ & @@NEVADA-RMSE@@ & @@NEVADA-R2@@ & @@NEVADA-NRMSE@@ \\
Tesla Texas & @@TEXAS-MAE@@ & @@TEXAS-MSE@@ & @@TEXAS-RMSE@@ & @@TEXAS-R2@@ & @@TEXAS-NRMSE@@ \\
\bottomrule
\end{tabular}
\end{table*}

\begin{figure*}[!t]
\centering
\IfFileExists{../figuras/mae_deepnpts_por_localidade.pdf}{%
  \includegraphics[width=0.78\textwidth]{mae_deepnpts_por_localidade.pdf}%
}{\fbox{\parbox[c][5.0cm][c]{0.75\textwidth}{\centering
Figura pendente: \texttt{mae\_deepnpts\_por\_localidade.pdf}}}}
\caption{MAE do DeepNPTS e do concorrente de menor MAE em cada localidade. A
escolha do concorrente local é descritiva e posterior ao teste; não constitui
uma regra prévia de seleção.}
\label{fig:deepnpts_local}
\end{figure*}

Pelo menor MAE local, os modelos tiveram a seguinte distribuição descritiva de
vitórias: \textbf{@@CONTAGEM-VITORIAS@@}. Especificamente,
\textbf{@@VENCEDORES-POR-LOCALIDADE@@}. Essa contagem não substitui o critério
primário da Eq.~\eqref{eq:macro} nem demonstra superioridade fora das origens
avaliadas.

\subsection{Previsões e incerteza ao longo do tempo}

A Fig.~\ref{fig:serie_byd} compara a referência modelada da NSRDB, o DeepNPTS e
\textbf{@@COMPARACAO-BYD@@} nas doze origens de Camaçari. O DeepNPTS apresentou
menor erro absoluto em \textbf{@@MESES-DEEPNPTS-MELHOR-BYD@@} dos doze meses e
MAE de \textbf{@@BYD-MAE@@ W/m$^2$}; o concorrente obteve
\textbf{@@BYD-MAE-COMPARACAO@@ W/m$^2$}. A melhor sobreposição visual ocorreu
em \textbf{@@COMENTARIO-SOBREPOSICAO-BYD@@}. A comparação mês a mês é auxiliar:
a conclusão quantitativa decorre das métricas sobre todas as origens.

\begin{figure}[!t]
\centering
\IfFileExists{../figuras/previsao_mensal_byd_camacari.pdf}{%
  \includegraphics[width=\columnwidth]{previsao_mensal_byd_camacari.pdf}%
}{\fbox{\parbox[c][3.8cm][c]{0.92\columnwidth}{\centering
Figura pendente: \texttt{previsao\_mensal\_byd\_camacari.pdf}}}}
\caption{GHI de referência da NSRDB e previsões consolidadas do DeepNPTS e de
\textbf{@@COMPARACAO-BYD@@} para Camaçari em 2024.}
\label{fig:serie_byd}
\end{figure}

A Tabela~\ref{tab:prob} reúne a avaliação probabilística. Para o intervalo
nominal de 90\%, o DeepNPTS apresentou cobertura de
\textbf{@@PICP-DEEPNPTS@@\%} e largura média de
\textbf{@@MPIW-DEEPNPTS@@ W/m$^2$}; o CRPS foi
\textbf{@@CRPS-DEEPNPTS@@ W/m$^2$}. O DeepAR alcançou, respectivamente,
\textbf{@@PICP-DEEPAR@@\%}, \textbf{@@MPIW-DEEPAR@@ W/m$^2$} e
\textbf{@@CRPS-DEEPAR@@ W/m$^2$}. Assim,
\textbf{@@INTERPRETACAO-PROBABILISTICA@@}.

\begin{table}[!t]
\centering
\caption{Métricas probabilísticas médias entre localidades.}
\label{tab:prob}
\small
\setlength{\tabcolsep}{5pt}
\begin{tabular}{lrrr}
\toprule
\textbf{Modelo} & \textbf{CRPS} & \textbf{PICP 90\%} & \textbf{MPIW} \\
 & \textbf{W/m$^2$} & \textbf{\%} & \textbf{W/m$^2$} \\
\midrule
DeepNPTS & @@CRPS-DEEPNPTS@@ & @@PICP-DEEPNPTS@@ & @@MPIW-DEEPNPTS@@ \\
DeepAR   & @@CRPS-DEEPAR@@ & @@PICP-DEEPAR@@ & @@MPIW-DEEPAR@@ \\
\bottomrule
\end{tabular}
\end{table}

\begin{figure}[!t]
\centering
\IfFileExists{../figuras/intervalo_deepnpts_byd_camacari.pdf}{%
  \includegraphics[width=\columnwidth]{intervalo_deepnpts_byd_camacari.pdf}%
}{\fbox{\parbox[c][3.8cm][c]{0.92\columnwidth}{\centering
Figura pendente: \texttt{intervalo\_deepnpts\_byd\_camacari.pdf}}}}
\caption{Mediana e intervalo preditivo central de 90\% do DeepNPTS para
Camaçari. A linha de referência corresponde ao produto modelado da NSRDB.}
\label{fig:intervalo_byd}
\end{figure}

\subsection{Limitações}

Há cinco limitações principais. Primeiro, quatro anos de alvos de treinamento
representam poucos ciclos sazonais. Segundo, a janela de 2024 já foi examinada
durante o desenvolvimento, o que restringe a interpretação a uma reanálise
retrospectiva. Terceiro, cinco sementes quantificam apenas uma parte da
variabilidade computacional. Quarto, a quantização impõe discretização aos
modelos aprendidos, enquanto as referências simples usam os valores físicos
exatos. Quinto, os dados da NSRDB são referência modelada, e as localidades das
fábricas não fornecem qualquer informação sobre geração, carga ou operação
industrial. Consequentemente, os resultados não devem ser extrapolados para
produção fotovoltaica ou para outras regiões sem nova validação.

\section{Conclusões}
\label{sec:conclusoes}

Este trabalho avaliou o DeepNPTS discreto oficial em um protocolo global de
previsão mensal de GHI. Dez séries relacionadas foram processadas com janelas
causais, transformação ajustada somente no pré-teste, cinco sementes e doze
origens retrospectivas de um mês à frente. A comparação incluiu três
referências simples e seis modelos aprendidos, além do DeepNPTS, e analisou
acurácia pontual e incerteza probabilística.

No protocolo estudado, \textbf{@@MODELO-MELHOR@@} obteve o menor macro-MAE,
\textbf{@@MAE-MELHOR@@ W/m$^2$}. O DeepNPTS alcançou
\textbf{@@MAE-DEEPNPTS@@ W/m$^2$}, ocupou a posição
\textbf{@@POSICAO-DEEPNPTS@@} e apresentou
\textbf{@@RESUMO-DEEPNPTS-VS-REFERENCIAS@@}. Na avaliação probabilística,
\textbf{@@RESUMO-DEEPNPTS-VS-DEEPAR@@}. Esses achados respondem à pergunta
experimental dentro da janela analisada, mas não estabelecem um modelo
universalmente superior.

Trabalhos futuros devem incorporar anos adicionais, reservar um período
prospectivo ainda não inspecionado, avaliar outras resoluções e horizontes e
examinar variáveis meteorológicas exógenas. Também é necessário estudar a
sensibilidade à quantização e ampliar o número de repetições antes de formular
conclusões externas sobre o DeepNPTS.

\begin{thebibliography}{99}

\bibitem{antonanzas2016}
J. Antonanzas et al., ``Review of photovoltaic power forecasting,''
\textit{Solar Energy}, vol. 136, pp. 78--111, 2016,
doi: 10.1016/j.solener.2016.06.069.

\bibitem{voyant2017}
C. Voyant et al., ``Machine learning methods for solar radiation forecasting:
A review,'' \textit{Renewable Energy}, vol. 105, pp. 569--582, 2017,
doi: 10.1016/j.renene.2016.12.095.

\bibitem{yang2018}
D. Yang et al., ``History and trends in solar irradiance and PV power
forecasting: A preliminary assessment and review using text mining,''
\textit{Solar Energy}, vol. 168, pp. 60--101, 2018,
doi: 10.1016/j.solener.2017.11.023.

\bibitem{duffie2013}
J. A. Duffie and W. A. Beckman, \textit{Solar Engineering of Thermal
Processes}, 4th ed. Hoboken, NJ, USA: Wiley, 2013.

\bibitem{sengupta2018}
M. Sengupta et al., ``The National Solar Radiation Data Base (NSRDB),''
\textit{Renewable and Sustainable Energy Reviews}, vol. 89, pp. 51--60, 2018,
doi: 10.1016/j.rser.2018.03.003.

\bibitem{voyant2022}
C. Voyant et al., ``Benchmarks for solar radiation time series forecasting,''
\textit{Renewable Energy}, vol. 191, pp. 747--762, 2022,
doi: 10.1016/j.renene.2022.04.065.

\bibitem{sobri2018}
S. Sobri, S. Koohi-Kamali, and N. A. Rahim, ``Solar photovoltaic generation
forecasting methods: A review,'' \textit{Energy Conversion and Management},
vol. 156, pp. 459--497, 2018, doi: 10.1016/j.enconman.2017.11.019.

\bibitem{chen2016}
T. Chen and C. Guestrin, ``XGBoost: A scalable tree boosting system,'' in
\textit{Proc. 22nd ACM SIGKDD Int. Conf. Knowledge Discovery and Data Mining},
2016, pp. 785--794, doi: 10.1145/2939672.2939785.

\bibitem{hornik1989}
K. Hornik, M. Stinchcombe, and H. White, ``Multilayer feedforward networks are
universal approximators,'' \textit{Neural Networks}, vol. 2, no. 5,
pp. 359--366, 1989, doi: 10.1016/0893-6080(89)90020-8.

\bibitem{elman1990}
J. L. Elman, ``Finding structure in time,'' \textit{Cognitive Science},
vol. 14, no. 2, pp. 179--211, 1990,
doi: 10.1207/s15516709cog1402\_1.

\bibitem{hochreiter1997}
S. Hochreiter and J. Schmidhuber, ``Long short-term memory,''
\textit{Neural Computation}, vol. 9, no. 8, pp. 1735--1780, 1997,
doi: 10.1162/neco.1997.9.8.1735.

\bibitem{chang2017}
S. Chang et al., ``Dilated recurrent neural networks,'' in
\textit{Advances in Neural Information Processing Systems}, vol. 30, 2017.

\bibitem{salinas2020}
D. Salinas, V. Flunkert, J. Gasthaus, and T. Januschowski, ``DeepAR:
Probabilistic forecasting with autoregressive recurrent networks,''
\textit{International Journal of Forecasting}, vol. 36, no. 3,
pp. 1181--1191, 2020, doi: 10.1016/j.ijforecast.2019.07.001.

\bibitem{rangapuram2023}
S. S. Rangapuram, J. Gasthaus, L. Stella, V. Flunkert, D. Salinas, Y. Wang,
and T. Januschowski, ``Deep non-parametric time series forecaster,''
\textit{arXiv preprint arXiv:2312.14657}, 2023,
doi: 10.48550/arXiv.2312.14657.

\bibitem{alexandrov2020}
A. Alexandrov et al., ``GluonTS: Probabilistic and neural time series modeling
in Python,'' \textit{Journal of Machine Learning Research}, vol. 21, no. 116,
pp. 1--6, 2020.

\bibitem{nlr2026}
National Laboratory of the Rockies, ``GOES Aggregated PSM v4 Download API,''
2026. [Online]. Available:
\url{https://developer.nlr.gov/docs/solar/nsrdb/nsrdb-GOES-aggregated-v4-0-0-download/}.
Accessed: Jul. 19, 2026.

\bibitem{hyndman2006}
R. J. Hyndman and A. B. Koehler, ``Another look at measures of forecast
accuracy,'' \textit{International Journal of Forecasting}, vol. 22, no. 4,
pp. 679--688, 2006, doi: 10.1016/j.ijforecast.2006.03.001.

\bibitem{gneiting2007}
T. Gneiting and A. E. Raftery, ``Strictly proper scoring rules, prediction, and
estimation,'' \textit{Journal of the American Statistical Association},
vol. 102, no. 477, pp. 359--378, 2007,
doi: 10.1198/016214506000001437.

\end{thebibliography}

\end{document}
